\paragraph{Phase Kickback with Controlled Phase Gates}\label{ex:phase-kickback-with-controlled-phase-gates}

In the previous section, we have seen the phase kickback using the CNOT gate. Since CNOT is a controlled-$X$ gate, the phase appeared because the target qubit was prepared in an eigenstate of $X$.

\highspace
Phase kickback, however, is \textbf{not limited to Pauli gates}. It can occur with \textbf{any controlled unitary operator}, provided that:
\begin{itemize}
    \item The \emph{target} qubit is an eigenstate of the unitary;
    \item The \emph{control} qubit is prepared in a superposition, usually $\ket{+}$.
\end{itemize}
A particularly simple example is obtained using a \textbf{controlled phase gate}.

\highspace
The \textbf{one-qubit} \definition{Phase Gate} is defined as:
\begin{equation}
    P\left(\phi\right) = \begin{pmatrix}
        1 & 0 \\
        0 & e^{i\phi}
    \end{pmatrix}
\end{equation}
Its action is straightforward:
\begin{equation*}
    P\left(\phi\right)\ket{0} = \ket{0}, \qquad P\left(\phi\right)\ket{1} = e^{i\phi}\ket{1}
\end{equation*}
Unlike the Pauli-$X$ gate, which changes the computational basis states:
\begin{equation*}
    X\ket{0} = \ket{1}, \qquad X\ket{1} = \ket{0}
\end{equation*}
The phase gate \textbf{does not change the basis state}. Instead, it leaves the computational basis unchanged while multiplying the state \ket{1} by the complex phase factor $e^{i\phi}$. For this reason, it is called a \emph{phase gate}.

\highspace
\textcolor{Green3}{\faIcon{book} \textbf{Eigenstates and Eigenvalues of the Phase Gate.}} The phase gate is a unitary operator, and its eigenstates are the computational basis states $\ket{0}$ and $\ket{1}$. The corresponding eigenvalues are simply the factors by which the eigenstates are multiplied:
\begin{equation*}
    P\left(\phi\right)\ket{0} = 1 \cdot \ket{0}, \qquad P\left(\phi\right)\ket{1} = e^{i\phi} \cdot \ket{1}
\end{equation*}
Thus, the eigenvalues of the phase gate are:
\begin{equation*}
    \lambda_0 = 1, \qquad \lambda_1 = e^{i\phi}
\end{equation*}

\highspace
The \definition{Controlled Phase Gate} is a two-qubit gate that applies the phase gate to the target qubit only when the control qubit is in the state $\ket{1}$. Its matrix representation is:
\begin{equation}
    CP\left(\phi\right) =
    \begin{pmatrix}
        I & 0 \\
        0 & P\left(\phi\right)
    \end{pmatrix} =
    \begin{pmatrix}
        1 & 0 & 0 & 0 \\
        0 & 1 & 0 & 0 \\
        0 & 0 & 1 & 0 \\
        0 & 0 & 0 & e^{i\phi}
    \end{pmatrix}
\end{equation}
As with every controlled gate:
\begin{itemize}
    \item If the control qubit is \ket{0}, the target qubit is left unchanged;
    \item If the control qubit is \ket{1}, the phase gate is applied to the target qubit.
\end{itemize}
The corresponding quantum circuit is:
\begin{center}
    \begin{quantikz}
        \lstick{\text{control}} & \ctrl{1} & \\
        \lstick{\text{target}} & \gate{P(\phi)} &
    \end{quantikz}
\end{center}

\highspace
\begin{examplebox}[: Target in \ket{0}]
    Since:
    \begin{equation*}
        P\left(\phi\right)\ket{0} = \ket{0}
    \end{equation*}
    The eigenvalue is $\lambda=1$. Therefore, if the target qubit is prepared in the state $\ket{0}$, no phase kickback occurs. The control qubit remains unchanged. Let's verify; prepare the state with the control qubit in $\ket{+}$ and the target qubit in $\ket{0}$:
    \begin{equation*}
        \ket{\psi} = \ket{+}\otimes\ket{0}
    \end{equation*}
    Applying the controlled phase gate:
    \begin{center}
        \begin{quantikz}
            \lstick{\ket{+}} & \ctrl{1} & \\
            \lstick{\ket{0}} & \gate{P(\phi)} &
        \end{quantikz}
    \end{center}
    \begin{align*}
        CP\left(\phi\right)\ket{\psi} &= CP\left(\phi\right)\left(\frac{1}{\sqrt{2}}\left(\ket{0}\otimes\ket{0} + \ket{1}\otimes\ket{0}\right)\right) \\
        &= \frac{1}{\sqrt{2}}\left(CP\left(\phi\right)\ket{00} + CP\left(\phi\right)\ket{10}\right) \\
        &= \frac{1}{\sqrt{2}}\left(\ket{00} + \ket{1} P(\phi)\ket{0}\right) \\
        &= \frac{1}{\sqrt{2}}\left(\ket{00} + \ket{1} \cdot 1 \cdot \ket{0}\right) \\
        &= \frac{1}{\sqrt{2}}\left(\ket{00} + \ket{10}\right) \\
        &= \frac{1}{\sqrt{2}}\left(\ket{0} + \ket{1}\right)\otimes\ket{0} \\
        &= \ket{\psi}
    \end{align*}
    As expected, the state remains unchanged, because the kicked-back phase is $\lambda=1$: $e^{i\cdot 0} = 1$. Exactly as for CNOT with the target in \ket{+}, the phase kickback is trivial.
\end{examplebox}

\highspace
\begin{examplebox}[: Target in \ket{1}]
    Since:
    \begin{equation*}
        P\left(\phi\right)\ket{1} = e^{i\phi}\ket{1}
    \end{equation*}
    The eigenvalue is $\lambda=e^{i\phi}$. Therefore, if the target qubit is prepared in the state $\ket{1}$, a phase kickback occurs. The control qubit is multiplied by the phase factor $e^{i\phi}$. Let's verify; prepare the state with the control qubit in $\ket{+}$ and the target qubit in $\ket{1}$:
    \begin{equation*}
        \ket{\psi} = \ket{+}\otimes\ket{1}
    \end{equation*}
    Applying the controlled phase gate:
    \begin{center}
        \begin{quantikz}
            \lstick{\ket{+}} & \ctrl{1} & \\
            \lstick{\ket{1}} & \gate{P(\phi)} &
        \end{quantikz}
    \end{center}
    \begin{align*}
        CP\left(\phi\right)\ket{\psi} &= CP\left(\phi\right)\left(\frac{1}{\sqrt{2}}\left(\ket{0}\otimes\ket{1} + \ket{1}\otimes\ket{1}\right)\right) \\
        &= \frac{1}{\sqrt{2}}\left(CP\left(\phi\right)\ket{01} + CP\left(\phi\right)\ket{11}\right) \\
        &= \frac{1}{\sqrt{2}}\left(\ket{01} + \ket{1} P(\phi)\ket{1}\right) \\
        &= \frac{1}{\sqrt{2}}\left(\ket{01} + e^{i\phi}\ket{11}\right) \\
        &= \frac{1}{\sqrt{2}}\left(\ket{0} + e^{i\phi}\ket{1}\right)\otimes\ket{1}
    \end{align*}
    The target qubit remains unchanged, while the phase has appeared on the \textbf{control qubit}. This is precisely the phenomenon of phase kickback. The kicked-back phase is $\lambda=e^{i\phi}$, as expected.
\end{examplebox}